\chapter{Semantic Entropy and Related Estimators}
\label{ch:semantic_entropy}

This chapter defines the method: discrete semantic entropy over sampled
answers. The aim is to keep the central idea of Farquhar et
al.~\cite{semantic_entropy_nature_2024}, which is to measure uncertainty over
meanings rather than over strings, while stating the procedure precisely enough that
the numbers in Chapter~\ref{ch:results} can be traced to specific steps of it.
The meaning-level framing follows Kuhn et
al.~\cite{kuhn_semantic_uncertainty_2023}.

Two commitments run through the chapter. Where the published description and
the released code~\cite{kossen_semantic_uncertainty_code_2024} disagree, the
code is what gets described, because a reader reproducing the result runs the
code and not the prose. The clustering step in Section~\ref{sec:clustering} is
where that matters. And the pipeline is kept in separable pieces, namely sampling,
grouping and scoring, so Chapter~\ref{ch:experimental_design} can vary one
while holding the others fixed. The last sections
(\ref{sec:related-estimators}--\ref{sec:other-estimators}) define the three
estimators compared against it.

\section{Entropy as a measure of uncertainty}
\label{sec:entropy-measure}

For a discrete distribution $p_1,\ldots,p_K$ over $K$ outcomes, the Shannon
entropy is
\[
  H(p) \;=\; -\sum_{k=1}^{K} p_k \log p_k .
\]
It is small when almost all the mass sits on one outcome and large when the
mass spreads over many. The maximum is $\log K$, at the uniform distribution.
The base of the logarithm changes the unit but not the ordering, so it does not
affect AUROC. This thesis uses the natural logarithm, so values are in nats.
Entropy is used because it is the quantity this literature defines and
reports~\cite{kuhn_semantic_uncertainty_2023,semantic_entropy_nature_2024}.
Nothing turns on the choice: AUROC depends only on the ranking, so any
increasing measure of spread would give similar results.

For semantic entropy the same formula applies, but the outcomes are not raw
strings. They are groups of strings taken to mean the same thing, so the
entropy measures how spread out the meanings are.

\subsection{Notation}

Table~\ref{tab:notation} lists the symbols used in this chapter.

\begin{table}[ht]
\thisfloatsetup{capposition=above}
\centering
\caption{Notation used for discrete semantic entropy.}
\label{tab:notation}
\begin{tabular}{@{}lp{8.2cm}@{}}
\hline
\textbf{Symbol} & \textbf{Meaning} \\
\hline
$x$ & prompt (question) supplied to the model \\
$M$ & number of sampled answers drawn from the model per prompt \\
$y_i$ & the $i$-th sampled answer, $i \in \{1,\ldots,M\}$ \\
$Y(x)$ & the multiset of all $M$ sampled answers for prompt $x$ \\
$\ell_i$ & length-normalized sequence log-likelihood of $y_i$ \\
$e(y_i, y_j)$ & NLI entailment decision: 1 if the entailment model labels
                        the pair $y_i \to y_j$ as \emph{entailment},
                        0 if it labels it \emph{neutral} or
                        \emph{contradiction} \\
$\sim_{\mathrm{sem}}$ & semantic equivalence relation used by the
                        clustering step (see Section~\ref{sec:clustering}) \\
$C(x)$ & the set of clusters $\{c_1,\ldots,c_K\}$ induced on $Y(x)$ \\
$K$ & number of distinct clusters, $1 \le K \le M$ \\
$n_k$ & number of sampled answers assigned to cluster $c_k$ \\
$\hat{p}_k$ & empirical cluster probability $n_k / M$ \\
$\widehat{H}_{\mathrm{sem}}(x)$ & discrete semantic entropy score \\
$\widehat{H}_{\mathrm{lex}}(x)$ & surface entropy score over normalized
                                  strings (the clustering-free comparator;
                                  called \emph{surface entropy} in
                                  Chapter~\ref{ch:results}) \\
AUROC & area under the receiver operating characteristic curve \\
AURAC & area under the rejection-accuracy curve \\
\hline
\end{tabular}
\end{table}

\section{The method, step by step}
\label{sec:method-steps}

The method has three computational steps. All three run on a fixed, trained
model and only read what it produces; none of them changes its weights.

\begin{enumerate}
  \item \textbf{Sampling.} For a prompt $x$, draw $M$ answers
        $y_1,\ldots,y_M$ under a fixed decoding configuration. The values used
        are in Section~\ref{sec:sampling-settings}.
  \item \textbf{Clustering.} Group the $M$ answers into $K$ clusters, putting
        two answers together if they are judged equivalent in the context
        of $x$.
  \item \textbf{Entropy scoring.} Compute each cluster's share
        $\hat{p}_k=n_k/M$ and the entropy
        $\widehat{H}_{\mathrm{sem}}(x)=-\sum_k \hat{p}_k\log\hat{p}_k$.
\end{enumerate}

Attaching a correctness label and using the score to decide what to answer
belongs to the \emph{evaluation}, not to the method, and
Chapter~\ref{ch:experimental_design} handles it separately. Keeping that line
is what makes the comparisons in Chapter~\ref{ch:results} possible: clustering
consumes stored samples, so the entailment model can be swapped without
regenerating anything, and correctness labelling consumes a finished answer, so
the grading rule can be swapped without touching the estimators.

Figure~\ref{fig:method-schematic} reproduces the original illustration. Panel
(a) makes the point on a short question: counting raw strings treats ``Paris''
and ``It's Paris'' as different answers and the entropy looks too high, while
grouping by meaning gives a low score. Panel (b) shows the long-form version,
where a biography is split into claims, each claim is re-asked as a question,
and the regenerated answers are grouped by meaning.

\begin{figure}[ht]
\centering
\includegraphics[width=\textwidth]{semantic_entropy_paper_fig1.png}
\caption{The semantic entropy method. (a)~On a short question, count\-ing
  distinct wordings makes the model look more un\-cer\-tain than it is; the
  source figure labels this ``mis\-lead\-ing\-ly high naive entropy''.
  Grouping by meaning first gives a low semantic entropy. (b)~The same idea
  applied to paragraph-length bi\-og\-ra\-phies, one claim at a time.
  Reproduced from Farquhar et al.~\cite{semantic_entropy_nature_2024}
  (\emph{Nature}, 2024), licensed under CC~BY~4.0.}
\label{fig:method-schematic}
\end{figure}

\section{Sampling setup}
\label{sec:sampling-setup}

For each prompt the model is sampled $M$ times under one decoding
configuration, giving $Y(x) = \{y_1,\ldots,y_M\}$. Two parameters control the
randomness. \emph{Temperature} $T$ divides the token logits before the softmax.
Low $T$ gives little variety, high $T$ gives more. \emph{Nucleus sampling}
(top-$p$) keeps only the smallest set of tokens whose probability adds up to at
least $p$.

Because the estimators read nothing but these samples, the temperature decides
how much signal there is to measure at all.
Section~\ref{sec:temperature-sensitivity} reports what happens across the
range, and Section~\ref{sec:sampling-settings} gives the values used.

\section{Clustering by equivalence}
\label{sec:clustering}

The main design choice is the equivalence relation. Two are used.

\paragraph{Normalized exact match.} The simplest relation:
$y_i \sim_{\mathrm{exact}} y_j$ when their normalized strings are identical.
Normalization is each dataset's standard answer-cleaning step. For the
question-answering datasets it lowercases, strips punctuation and the articles
\emph{a\,/\,an\,/\,the}, and collapses whitespace, following the TriviaQA
evaluation protocol, so ``The Eiffel Tower!'' and ``eiffel tower'' become the
same string. For the arithmetic dataset it extracts the final number, so ``The
answer is 7 apples'' becomes ``7''. This is a surface grouping: it collapses
spelling and formatting and nothing else. The entropy over these groups is
called \emph{surface entropy}, and it is the clustering-free comparator that
RQ2 turns on. A cruder companion, \emph{naive sample entropy}, counts the raw
strings without normalizing first.

\paragraph{Bidirectional entailment.} Instead of string identity, this rule
asks whether two answers entail each other. For $y_i$ and $y_j$ it asks
whether $y_i$ entails $y_j$ \emph{and} whether $y_j$ entails $y_i$. If both
hold, they are the same semantic answer:
\[
  y_i \sim_{\mathrm{sem}} y_j
  \quad\Longleftrightarrow\quad
  e(y_i,y_j)=1 \ \text{and}\ e(y_j,y_i)=1.
\]
Farquhar et al.~\cite{semantic_entropy_nature_2024} define equivalence
\emph{in the context of the question}, and give the reason: ``Paris'' does not
entail ``The capital of France is Paris'' on its own, but in the context of the
question it does. The short-answer results here follow that definition. Every
entailment query in the canonical configuration is prefixed with the question,
and the LLM backend is given the question alongside both answers for the same
reason. The one exception is the whole-paragraph biography protocol, which is
answer-only, because a generated paragraph is not a response to a single
question.

Two details of the input are easy to assume wrongly. The entailment step
consumes the \emph{raw} generations, not the normalized strings above:
normalization belongs to the comparator, and applying it first would confound
the two. And the question is joined to each answer as plain text, so what the
entailment model reads is a pair of question-prefixed strings.

The rule approximates equivalence rather than defining it, because judgments
from a small cross-encoder can be noisy. It may split answers that differ only
in harmless detail, and merge answers that hide real differences. Four backends
compute $e(\cdot,\cdot)$:

\begin{itemize}
  \item \texttt{microsoft/deberta-v2-xlarge-mnli}, the model the released code
    loads by default~\cite{he_deberta_2021,williams_multinli_2018}. This is the
    \emph{canonical} backend: results reported without a named backend use it,
    because it is what a reader running the published code would get.
  \item \texttt{cross-encoder/nli-deberta-v3-large}, a current
    cross-encoder~\cite{sentence_transformers_nli_deberta_v3_large_2022,%
    he_debertav3_2021}.
  \item \texttt{cross-encoder/nli-deberta-v3-base}, a smaller
    cross-encoder~\cite{sentence_transformers_nli_deberta_v3_base_2022}, used
    on the long-form arm only.
  \item \texttt{Qwen/Qwen2.5-72B-Instruct}~\cite{yang_qwen25_2024}, an open
    stand-in for the proprietary judge Farquhar et al.\ use on sentence-length
    answers. It is prompted with the entailment instruction taken verbatim from
    the released code.
\end{itemize}

Every backend is applied afterwards to one stored set of generations, so the
entailment model is the only thing that differs between them and any change in
the scores comes from it alone. Section~\ref{sec:clustering-sensitivity}
quantifies how much that choice matters, and it turns out to matter more than
the difference between the estimators.

\subsection{The grouping procedure as implemented}
\label{sec:anchor-scan}

The grouping follows the released
code~\cite{kossen_semantic_uncertainty_code_2024} rather than a tidier
equivalent. Answers are scanned in generation order. The first answer without a
cluster identifier becomes an \emph{anchor} and gets a new identifier, and
every later answer mutually entailed with that anchor gets the same one. The
scan then moves to the next unassigned answer.

Two properties matter, and neither is incidental.

First, the assignment in the inner loop is unconditional. The released routine
guards only the \emph{outer} loop, which decides whether an answer becomes an
anchor; the inner assignment carries no such test. An answer a previous anchor
already claimed is therefore \emph{reassigned} when a later anchor also matches
it, so its final cluster is decided by the last anchor that matches rather than
the first.

Second, the procedure computes no transitive closure. If A is mutually entailed
with B, and B with C, but A and C are not, then A, B and C need not end up in
one group; whether they do depends on which the scan reaches first. The result
is \emph{order-dependent}: permuting the samples can change the number of
clusters, and so the entropy, without changing a single entailment judgment.

An earlier version of this work grouped answers by the connected components of
the mutual-entailment graph, using union-find. That is the natural reading of
``group answers that mean the same thing'', it is order-independent, and it
merges transitively. It is also not what the code does, and the two procedures
give different cluster counts on the same entailment matrix whenever the
relation is not transitive, which for a noisy neural model is common.
The point of this thesis is to measure the method as published, so the anchor
scan is used throughout and its order dependence is carried as a property of
the method under study. Algorithm~\ref{alg:clustering} states it
exactly.

\begin{algorithm}
\caption{Bidirectional-entailment clustering by order-dependent anchor scan,
  reproducing \texttt{get\_semantic\_ids} in the released
  code~\cite{kossen_semantic_uncertainty_code_2024}.}
\label{alg:clustering}
\KwIn{answers $y_1,\ldots,y_M$ in generation order; entailment model
  $\mathcal{M}(\cdot,\cdot)$ returning entailment, neutral, or
  contradiction (queries are prefixed with the question in the
  question-conditioned configuration)}
\KwOut{one cluster identifier per answer}
$c_i \gets \textsc{unassigned}$ for $i=1,\ldots,M$\;
$k \gets 0$\tcp*{next cluster identifier}
\For{$i \gets 1$ \KwTo $M$}{
  \If{$c_i$ is \textsc{unassigned}}{
    $c_i \gets k$\tcp*{$y_i$ becomes an anchor}
    \For{$j \gets i+1$ \KwTo $M$}{
      $a \gets \mathcal{M}(y_i, y_j)$\tcp*{does $y_i$ entail $y_j$?}
      $b \gets \mathcal{M}(y_j, y_i)$\tcp*{and $y_j$ entail $y_i$?}
      \If{$a$ is entailment \textbf{and} $b$ is entailment}{
        $c_j \gets k$\tcp*{assigned even if $y_j$ already had a cluster}
      }
    }
    $k \gets k + 1$\;
  }
}
\Return $(c_1,\ldots,c_M)$\;
\end{algorithm}

\section{Discrete semantic entropy}
\label{sec:discrete-se}

After clustering, let $n_k$ be the number of answers in cluster $c_k$. Its
empirical probability is $\hat{p}_k = n_k/M$, and the discrete semantic entropy
is
\[
  \widehat{H}_{\mathrm{sem}}(x)
  \;=\; -\sum_{k=1}^{K} \hat{p}_k \log \hat{p}_k .
\]
A low score means the samples mostly agree on one meaning; a high score means
they spread across several. The maximum is $\log K \le \log M$. The experiments
use $M=10$, so at most $\ln 10 \approx 2.30$; the worked example below uses
$M=5$, so at most $\ln 5 \approx 1.61$.

\paragraph{The estimator family.}\label{sec:estimator-family}
The grid reports this count-based estimator, the probability-weighted variant
of Farquhar et al.~\cite{semantic_entropy_nature_2024,%
kossen_semantic_uncertainty_code_2024}, and two clustering-free scores. Two
axes separate them: whether answers are grouped by meaning, and whether the
score reads sequence probabilities or only counts.

\begin{center}
\small
\resizebox{\linewidth}{!}{%
\begin{tabular}{@{}lll@{}}
\toprule
 & Count-based & Probability-weighted \\
\midrule
No meaning judgment & \makecell[l]{naive sample entropy (raw strings),\\
  surface entropy (normalized strings)} &
  \makecell[l]{naive entropy\\(the paper's naive baseline)} \\
Meaning judgment & \makecell[l]{discrete semantic entropy\\
  (the paper's discrete approximation)} &
  \makecell[l]{full semantic entropy\\(the paper's primary estimator)} \\
\bottomrule
\end{tabular}}
\end{center}

Two of the four ignore meaning and read only the strings: \emph{naive sample
entropy} counts the raw strings, \emph{surface entropy} counts them after
normalization. A third, the paper's \emph{naive entropy}, ignores meaning
differently: it does no grouping at all and is the mean negative
length-normalized sequence log-likelihood over the $M$ samples, so it reads how
confident the model was rather than how often an answer recurred. The two
semantic scores group first: \emph{discrete semantic entropy} counts the
groups, \emph{full semantic entropy} weights them by the summed likelihood of
their members. When the samples are equally likely the full version reduces
exactly to the discrete one, which is a useful check on the implementation.

The distinction between naive entropy and the two string scores is what RQ2
rests on, and it is easy to lose. Naive entropy differs from discrete semantic
entropy in \emph{two} ways: it reads likelihoods rather than counts, and it does
no grouping. Surface entropy differs in \emph{one}: it counts exactly as the
discrete estimator does, but groups by string identity rather than by meaning.
So a comparison against naive entropy cannot separate the contribution of the
grouping from the contribution of counting. A comparison against surface
entropy can.

Algorithm~\ref{alg:scoring} gives the procedure. Note the ordering: clustering
consumes the raw answers, and normalization is applied only on the branch that
produces the comparator.

\begin{algorithm}
\caption{Semantic-entropy scoring from sampled answers, with the
  clustering-free comparator computed from the same samples.}
\label{alg:scoring}
\KwIn{raw sampled answers $y_1,\ldots,y_M$ for prompt $x$; optional
  length-normalized sequence log-likelihoods $\ell_1,\ldots,\ell_M$;
  entailment model $\mathcal{M}$}
\KwOut{discrete semantic entropy $\widehat{H}_{\mathrm{sem}}$, surface
  entropy $\widehat{H}_{\mathrm{lex}}$, and, when log-likelihoods are
  available, full semantic entropy $\widehat{H}^{\mathrm{p}}_{\mathrm{sem}}$}
$(c_1,\ldots,c_M) \gets \textsc{cluster}(y_1,\ldots,y_M;\mathcal{M})$\tcp*{Algorithm~\ref{alg:clustering}, on raw answers}
let $K$ be the number of distinct clusters\;
\For{$k \gets 1$ \KwTo $K$}{
  $n_k \gets |\{\,i : c_i = k\,\}|$\;
  $\hat{p}_k \gets n_k / M$\;
}
$\widehat{H}_{\mathrm{sem}} \gets -\sum_{k=1}^{K}\hat{p}_k \log \hat{p}_k$\tcp*{discrete semantic entropy}
\tcp{comparator: normalize, then count distinct strings}
$u_i \gets \textsc{normalize}(y_i)$ for $i=1,\ldots,M$\;
$\hat{q}_s \gets |\{\,i : u_i = s\,\}| / M$ for each distinct string $s$\;
$\widehat{H}_{\mathrm{lex}} \gets -\sum_{s}\hat{q}_s \log \hat{q}_s$\tcp*{surface entropy}
\If{sequence log-likelihoods are available}{
  $\log P_k \gets \operatorname{logsumexp}_{\,i:\,c_i=k}\ \ell_i$ for each $k$\;
  $Z \gets \operatorname{logsumexp}_{\,k}\ \log P_k$\tcp*{renormalize cluster masses}
  $q_k \gets \exp(\log P_k - Z)$ for each $k$\;
  $\widehat{H}^{\mathrm{p}}_{\mathrm{sem}} \gets -\sum_{k=1}^{K} q_k \log q_k$\tcp*{full semantic entropy}
}
\Return $\widehat{H}_{\mathrm{sem}}$, $\widehat{H}_{\mathrm{lex}}$ (and
  $\widehat{H}^{\mathrm{p}}_{\mathrm{sem}}$ if computed)\;
\end{algorithm}

\section{Worked example}
\label{sec:worked_example}

The following example makes the pipeline concrete. Consider the factual
question $x = $ ``What is the capital of Australia?''. The language model is
sampled $M=5$ times under one decoding configuration, drawing a fresh
completion each time. That yields five different but plausible answers:

\begin{center}
\begin{tabular}{cl}
\hline
Index & Sampled answer $y_i$ \\
\hline
1 & \emph{Canberra} \\
2 & \emph{Canberra is the capital city of Australia} \\
3 & \emph{Sydney} \\
4 & \emph{Melbourne} \\
5 & \emph{Canberra} \\
\hline
\end{tabular}
\end{center}

The example is intentionally small and intentionally a case in which the model
is mostly right but somewhat uncertain, which makes it a useful check on the
pipeline. The decoding settings of the recorded experiments are given in
Section~\ref{sec:sampling-settings}.

\paragraph{Step 1: Bidirectional entailment.}
The entailment model is applied to each ordered pair $(y_i, y_j)$.
Table~\ref{tab:entailment_example} shows the resulting binary decisions; a
cell $(i,j)$ equals 1 if the row answer entails the column answer.

\begin{table}[ht]
\thisfloatsetup{capposition=above}
\centering
\caption{Pairwise entailment matrix for the worked example.
         A cell $(i,j)=1$ means $y_i$ entails $y_j$.}
\label{tab:entailment_example}
\begin{tabular}{c|ccccc}
       & $y_1$ & $y_2$ & $y_3$ & $y_4$ & $y_5$ \\\hline
$y_1$ &  --   &   1   &   0   &   0   &   1   \\
$y_2$ &   1   &  --   &   0   &   0   &   1   \\
$y_3$ &   0   &   0   &  --   &   0   &   0   \\
$y_4$ &   0   &   0   &   0   &  --   &   0   \\
$y_5$ &   1   &   1   &   0   &   0   &  --   \\
\end{tabular}
\end{table}

\paragraph{Step 2: Semantic clustering.}
Two answers belong to the same cluster when both off-diagonal cells are 1.
Reading from Table~\ref{tab:entailment_example}: $y_1 \sim y_2$,
$y_1 \sim y_5$ and $y_2 \sim y_5$, while $y_3$ and $y_4$ each stand alone.

It is worth tracing the anchor scan of Algorithm~\ref{alg:clustering} rather
than assuming the answer, since the scan and a connected-components grouping
do not agree in general. Answer $y_1$ is unassigned, so it anchors cluster~0;
the inner loop finds $y_2$ and $y_5$ mutually entailed with it and assigns
both to cluster~0, and rejects $y_3$ and $y_4$. The scan moves on: $y_2$ is
already assigned, $y_3$ is not and anchors cluster~1, matching nothing;
$y_4$ anchors cluster~2, matching nothing; $y_5$ is already assigned. The
entailment relation here happens to be transitive on
$\{y_1,y_2,y_5\}$, so no reassignment occurs and the scan returns the same
partition a connected-components grouping would. This yields $K=3$ clusters:
\[
  c_1 = \{y_1, y_2, y_5\},\quad
  c_2 = \{y_3\},\quad
  c_3 = \{y_4\}.
\]
Section~\ref{sec:order_dependence_example} shows a case in which that
agreement fails.

\paragraph{Step 3: Empirical probabilities.}
\[
  \hat{p}_1 = \tfrac{3}{5} = 0.60,\quad
  \hat{p}_2 = \tfrac{1}{5} = 0.20,\quad
  \hat{p}_3 = \tfrac{1}{5} = 0.20.
\]

\paragraph{Step 4: Discrete semantic entropy.}
\begin{align*}
  \widehat{H}_{\mathrm{sem}}(x)
  &= -\bigl(0.60\ln 0.60 + 0.20\ln 0.20 + 0.20\ln 0.20\bigr) \\
  &= -(-0.306 - 0.322 - 0.322) \\
  &\approx 0.95 \text{ nats.}
\end{align*}

\paragraph{The clustering-free comparator.}
This is not a step of the semantic-entropy method but the comparator RQ2
turns on, computed from the same five samples. It treats each distinct
normalized string as its own outcome. With strings
$\{$\emph{Canberra}$\times 2$, \emph{Canberra is the capital city of
Australia}$\times 1$, \emph{Sydney}$\times 1$, \emph{Melbourne}$\times 1\}$,
the surface entropy is
\begin{align*}
  \widehat{H}_{\mathrm{lex}}(x)
  &= -\bigl(0.40\ln 0.40 + 0.20\ln 0.20 + 0.20\ln 0.20 + 0.20\ln 0.20\bigr) \\
  &\approx 1.33 \text{ nats.}
\end{align*}

In this example the count over normalized strings and the count over raw
strings coincide, because no two of the five samples differ only in case or
punctuation. Chapter~\ref{ch:results} keeps them apart as \emph{surface
entropy} and \emph{naive sample entropy}, since on the recorded data they do
not always agree.

The surface score is higher (1.33 against 0.95) because it treats the
paraphrase ``Canberra is the capital city of Australia'' as a different answer
from ``Canberra'', inflating the apparent uncertainty. After semantic
clustering the correct answer class accounts for 60\,\% of the samples, and
the entropy is correspondingly lower. This is the mechanism the method is
named for, and the example is constructed to display it. Whether it produces a
measurable advantage on real data, once both scores are computed over tens of
thousands of records, is the question Chapter~\ref{ch:results} settles;
the example establishes what is being claimed, not that the claim holds.

\subsection{When the grouping depends on the order}
\label{sec:order_dependence_example}

The example above is well behaved: its entailment relation happens to be
transitive, so the anchor scan of Algorithm~\ref{alg:clustering} and a
connected-components grouping return the same three clusters. That is not
guaranteed, and when it fails the consequences are larger than they might
appear. The following case shows why.

Take four sampled answers to the same question:

\begin{center}
\begin{tabular}{cl}
\hline
Index & Sampled answer \\
\hline
$a$ & \emph{Canberra} \\
$b$ & \emph{Canberra, the capital} \\
$c$ & \emph{The capital city} \\
$d$ & \emph{Sydney} \\
\hline
\end{tabular}
\end{center}

Suppose the entailment model judges $a \leftrightarrow b$ mutually entailed
and $b \leftrightarrow c$ mutually entailed, but \emph{not}
$a \leftrightarrow c$: ``Canberra'' on its own does not entail ``the capital
city'', and vice versa, without the question supplying the link that $b$
states explicitly. This is not a contrived failure. It is the ordinary
behavior of a relation that is symmetric by construction but not transitive,
and a neural entailment model has no mechanism enforcing transitivity.

Grouping by connected components merges all three through the bridge $b$,
giving $\{a,b,c\},\{d\}$ regardless of the order the samples arrived in. The
anchor scan does not, and what it returns depends on that order:

\begin{center}
\begin{tabular}{lll}
\hline
Order presented & Resulting clusters & $\widehat{H}_{\mathrm{sem}}$ \\
\hline
$a,b,c,d$ & $\{a,b\},\{c\},\{d\}$      & $1.04$ nats \\
$b,a,c,d$ & $\{a,b,c\},\{d\}$          & $0.56$ nats \\
$a,c,b,d$ & $\{a\},\{b,c\},\{d\}$      & $1.04$ nats \\
\hline
\end{tabular}
\end{center}

Three permutations of the same four answers give the same entailment judgments
throughout, but two different partitions. One of them assigns nearly twice the
entropy of the other.

The third row is the more striking case, and it is the reassignment rule at
work. Answer $a$ anchors first and claims $b$. The scan then reaches $c$,
which is still unassigned, so $c$ becomes an anchor in its own right and
claims $b$ a second time and overwrites $b$'s existing identifier. Answer $a$ is
left in a cluster by itself, even though it is mutually entailed with an answer
it originally grouped with.

Two points follow. First, the reported entropy is a property of the generation
order as well as of the entailment judgments, which is worth knowing when the
estimator is described as measuring ``the number of distinct meanings''.
Second, the effect is invisible in aggregate: the partition is well defined
and the entropy is finite, so nothing downstream signals that a different
ordering would have produced a different number.

Whether this matters empirically is a separate question from whether it
occurs, and it is settled by measurement rather than argument. The
disagreement between entailment backends reported in
Section~\ref{sec:clustering-sensitivity} bounds how much the clustering step
moves the results, and order effects are one component of that. This example
establishes only that the estimator as published is order-dependent. A reader
who assumes otherwise, as the natural reading of ``cluster by meaning''
invites, would be mistaken.

\section{Method contract}
\label{sec:method-contract}

The algorithm itself needs little. It takes a prompt $x$ with $M$ sampled
answers drawn under a fixed decoding configuration, and optionally their
length-normalized sequence log-likelihoods. It groups the \emph{raw} answers
with the configured equivalence relation, each entailment query prefixed with
$x$ in the question-conditioned configuration, applying no normalization first.
It returns the cluster assignments and the discrete semantic entropy over
cluster frequencies, plus the full semantic entropy when log-likelihoods are
supplied. From the same samples it also applies the dataset's answer-cleaning
and returns the surface entropy over distinct normalized strings.

Correctness labelling and evaluation consume those outputs but are not part of
the method. That separation is what lets
Chapter~\ref{ch:experimental_design} apply four entailment backends and three
correctness rules to one stored set of generations.

\section{Related estimators evaluated alongside}
\label{sec:related-estimators}

The remaining estimators are published baselines, defined here only as far as
Chapter~\ref{ch:results} needs. One of them, the Semantic Entropy Probe,
addresses the main method's cost: discrete semantic entropy needs $M$ samples
per question and up to $M^2$ pairwise entailment checks. The other two, P(True)
and embedding regression, are simple and widely used.

All three read a single answer per question rather than the ten the entropy
estimators use. That answer is called the \emph{most-likely answer} throughout,
and it is also what correctness is scored on. It is not a greedy decode:
following the released code, it is a sample drawn at temperature $0.1$, under a
token budget four times larger than the samples get.
Section~\ref{sec:sampling-settings} gives the details.

\subsection{Semantic Entropy Probes}
\label{sec:sep-method}

A Semantic Entropy Probe (SEP)~\cite{kossen_semantic_entropy_probes_2024}
attacks the cost problem from the other end: it reads uncertainty off the
model's internal state, so at test time one forward pass is enough, with no
multi-sampling and no entailment comparisons.

For each training question, the most-likely answer supplies a hidden state,
read at the last generated token as fixed by the stop sequence. Separately, ten
sampled responses are grouped by meaning to compute discrete semantic entropy.
That entropy is turned into a high/low label, and a logistic-regression
classifier is trained to predict it from the hidden state. Correctness labels
are not used; the accuracy probe below is the variant trained on those. At test
time, generate one answer, read the same hidden state, apply the classifier,
and use its output as the uncertainty score.

The probe is trained separately per dataset, model and seed on its own prompts
and scored on a disjoint set, so its numbers are a genuine held-out test. The
feature choices are fixed in advance rather than tuned on the evaluation set.
One detail follows the code rather than the natural reading: the hidden state is
taken at position $n_{\mathrm{generated}}-1$, where $n_{\mathrm{generated}}$
comes from where the stop sequence cut the answer, and no end-of-sequence token
has to be present. The probe is far cheaper than full semantic entropy, but
limited to what is linearly readable from one hidden state.

\subsection{P(True) and embedding regression}
\label{sec:other-estimators}

\textbf{P(True)} asks the model to grade itself. It is shown the question, a
set of brainstormed alternative answers, and its own most-likely answer, and is
asked whether that answer is true, with the options presented as letters; a
twenty-shot prefix from the training split comes first. What the released code
records is the raw log-probability the model puts on the token for the ``true''
option, not a probability renormalized over the two options, which is what the
published description suggests. This thesis follows the code, and
Section~\ref{sec:deviation-register} records the difference. The score is
oriented like every other score here, larger meaning riskier. P(True) needs no
extra samples and no probe training, which makes it a natural cheap reference.

\textbf{Embedding regression} is a close cousin of the SEP probe: the same kind
of linear model on the same hidden states of the same most-likely answers, but
trained on \emph{judged correctness} instead of a semantic-entropy threshold.
It is reported in two forms. \emph{In-distribution} means trained and tested on
the same task. \emph{Out-of-distribution} means trained on the pooled prompts of
the other two short-answer tasks and scored on the held-out one. Together they
measure how well such a probe transfers between datasets.
Chapter~\ref{ch:results} calls it the \emph{accuracy probe}, since its training
target is what distinguishes it from SEP.
